\documentclass[12pt,a4paper]{article}

\usepackage{fontspec}
\usepackage{polyglossia}
\setdefaultlanguage{english}
\setmainfont{CMU Serif}
\setsansfont{CMU Sans Serif}
\setmonofont{CMU Typewriter Text}
\usepackage{geometry}
\geometry{top=2cm, bottom=2cm, left=2.5cm, right=2.5cm}
\usepackage{amsmath,amssymb,amsthm}
\usepackage{booktabs,longtable,float,xcolor}
\usepackage{hyperref}
\hypersetup{colorlinks, urlcolor=blue, citecolor=blue, linkcolor=black}
\usepackage{enumitem}
\sloppy
\tolerance=1000
\emergencystretch=3em
\numberwithin{equation}{section}

\title{\textbf{Model MEML~75:}\\[0.3cm]
\large A phenomenological parametrization of the Standard Model\\
fermion sector from $T^2/\mathbb{Z}_3$ geometry\\
with axion-like (ALP) dark matter}
\author{D.~V.~Yurov}
\date{2026}

\begin{document}
\maketitle

\begin{abstract}
\noindent
We present model MEML~75, which parametrizes the fermion sector of the
Standard Model through the geometry of the six-dimensional
orbifold $T^2/\mathbb{Z}_3$ with three independent tori.
The model fits \textbf{13 fermion observables} (9 quarks $+$ 4 leptons)
with \textbf{16 free parameters} ($N_{\rm dof} = 13-16 = -3$).

\textbf{The main testable prediction:}
axion-like dark matter with
$f_a = 1.9801\times10^{11}$~GeV,
$m_a = 28.407$~$\mu$eV, $\nu_a = 6.8688$~GHz, and ALP-photon
coupling
$g_{a\gamma\gamma} = +4.38\times10^{-15}$~GeV$^{-1}$ (DFSZ)
or $-1.13\times10^{-14}$~GeV$^{-1}$ (KSVZ), testable
at ADMX/HAYSTAC (0.5--10~GHz).

\textbf{Quark sector.} A global fit with loss function
$\chi^2 = \sum_i ((O_i-O_i^{\rm target})/\sigma_i)^2$ yields
$\chi^2 = 6.195\times10^{-20}$ at
$\tau_u = -0.1577944085 + 2.3851782590\,i$,
$\tau_d = +0.3693905617 + 1.5410496057\,i$.

\textbf{Lepton sector.} Using a third torus \(\tau_l\), a combination of $A_4$ modular forms $Y^{(2)}$ and $Y^{(4)}$ yields
$\chi^2 = 1.0065\times10^{-16}$ for the PMNS observables at
$\tau_l = 0.385257 + 0.958275\,i$. The model predicts
$\theta_{12}=33.41^\circ$, $\theta_{13}=8.54^\circ$, $\theta_{23}=49.10^\circ$, and $J_{cp}=-0.018$.

\textbf{Verdict:} the model is \emph{phenomenological}. It interpolates fermion observations with high precision.
\end{abstract}

\newpage

\section{Introduction}

Modular models of fermion masses~\cite{Feruglio2017,Kobayashi2019}
build the mass hierarchy and mixing from compactification geometry.
In model MEML~75, quarks and leptons are localized on three
independent tori $T^2/\mathbb{Z}_3$ with moduli $\tau_u$, $\tau_d$, and \(\tau_l\). The mass matrices are built from $A_4$ modular forms of weights $(2,4,6)$.

\section{Geometry of $T^2/\mathbb{Z}_3$}
The orbifold $T^2/\mathbb{Z}_3$ is defined by:
\begin{equation}
z = x + \tau y, \quad ds^2 = dx^2 + dx\,dy + dy^2,
\end{equation}
where $\tau = e^{i\pi/3} = \frac12 + i\frac{\sqrt3}{2}$. The $Z_3$-action $z \to \omega z$ (with $\omega = e^{2\pi i/3}$) has three fixed points, providing a topological explanation for three generations.

\section{Bulk coupling $\alpha_G$ (postulate)}
\begin{equation}
\alpha_G = \frac{2}{3\pi} = 0.2122065908.
\end{equation}
Geometric motivation: $2/3$ is the volume fraction of the fundamental domain of the $\mathbb{Z}_3$-orbifold; $1/\pi$ is the normalization of the theta-function period.

\section{Suppression factor $\varepsilon$ (ansatz)}
\begin{equation}
\varepsilon = \left(\frac{\alpha_G}{4\pi}\right)^2 = \frac{1}{36\pi^4} = 2.851662\times10^{-4}.
\end{equation}

\section{Conditionally derived quantities}
\begin{align}
m_{KK} &= \frac{M_{Pl}^{\rm red}}{\sqrt{\mathcal{I}}\,e^{\pi^2}}
= 2.0838\times10^{14}~\text{GeV}, \\
M_6 &= \left(3 (M_{Pl}^{\rm red})^2 m_{KK}^2\right)^{1/4}
= 2.9646\times10^{16}~\text{GeV},
\end{align}
at $\mathcal{I} = 0.3653$.

\section{Axion-like (ALP) sector}
\subsection{ALP parameters}
\begin{align}
f_a &= M_{Pl}^{\rm red}\,\varepsilon^2 = 1.9801\times10^{11}~\text{GeV}, \\
m_a &= \frac{(\chi^{1/4})^2}{f_a} = 28.407~\mu\text{eV}, \\
\nu_a &= \frac{m_a c^2}{h} = 6.8688~\text{GHz}.
\end{align}
Here $\chi^{1/4} \approx 75.0$~MeV is the fourth root of the topological susceptibility of QCD. We emphasize that no explicit \(U(1)_{PQ}\) symmetry is introduced; the object is an ALP.

\section{Quark sector (11 parameters)}
\subsection{Matrix structure}
\begin{equation}
M_u = \begin{pmatrix}
a_u Y_1^{(2)} & b_u Y_3^{(4)} & c_u Y_2^{(6)} \\
a_u Y_3^{(2)} & b_u Y_2^{(4)} & c_u Y_1^{(6)} \\
a_u Y_2^{(2)} & b_u Y_1^{(4)} & c_u Y_3^{(6)}
\end{pmatrix}, \;
M_d = \begin{pmatrix}
a_1 Y_1^{(6)} & b Y_3^{(2)} & c Y_2^{(2)} \\
a_1 Y_3^{(6)} & b Y_2^{(2)} & c Y_1^{(2)} \\
a_2 Y_2^{(2)} & b Y_1^{(2)} & c Y_3^{(4)}
\end{pmatrix}.
\end{equation}

\subsection{Results}
Independent optimization yields $\chi^2 = 6.195\times10^{-20}$ at the moduli given in the abstract.

\section{Lepton sector (5 parameters)}
\label{sec:lepton}
The charged lepton masses are taken diagonal (in the flavor basis), while the neutrino mass matrix \(M_\nu\) is constructed from the combination of weight-2 and weight-4 modular forms.

\subsection{Matrix structure}
\begin{equation}
M_\nu = g_1 M_{3_S}(Y^{(2)}) + g_2 M_{3_S}(Y^{(4)}) + d \cdot \mathbb{I},
\end{equation}
where the triplet contraction \(M_{3_S}(Y)\) is given by:
\begin{equation}
M_{3_S}(Y) = \begin{pmatrix}
2 Y_1 & -Y_3 & -Y_2 \\
-Y_3 & 2 Y_2 & -Y_1 \\
-Y_2 & -Y_1 & 2 Y_3
\end{pmatrix}.
\end{equation}
The diagonal term \(d \cdot \mathbb{I}\) breaks the \(\mu-\tau\) symmetry.

\subsection{Results}
Independent optimization yields:
\begin{equation}
\boxed{\;\chi^2 = 1.0065\times10^{-16}
\quad\text{at}\quad
\tau_l = 0.385257 + 0.958275\,i,\;
g_1 = 9.9009,\; g_2 = 9.6201,\; d = -7.0956.\;}
\end{equation}

\begin{table}[H]
\centering
\caption{PMNS observables (NuFIT 5.2).}
\begin{tabular}{lccc}
\toprule
Observable & Predicted & Target & Deviation \\
\midrule
$\theta_{12}$ (solar)    & $33.41^\circ$ & $33.41^\circ \pm 0.5^\circ$ & $\sim 10^{-9}$ \\
$\theta_{13}$ (reactor)  & $8.54^\circ$  & $8.54^\circ \pm 0.1^\circ$  & $\sim 10^{-9}$ \\
$\theta_{23}$ (atmos.)   & $49.10^\circ$ & $49.1^\circ \pm 0.5^\circ$  & $\sim 10^{-9}$ \\
$J_{cp}$ (CP phase)      & $-0.0180$     & $-0.018 \pm 0.005$          & $\sim 10^{-9}$ \\
\bottomrule
\end{tabular}
\end{table}

\section{Full balance}
\label{sec:balance}

\begin{table}[H]
\centering
\caption{Free parameters (16).}
\begin{tabular}{llc}
\toprule
Category & Parameters & Count \\
\midrule
Quark moduli & $\tau_u, \tau_d$ (complex) & 4 \\
Quark coefficients & $a_u, b_u, c_u, a_1, a_2, b, c$ & 7 \\
Lepton modulus & \(\tau_l\) (complex) & 2 \\
Lepton coefficients & $g_1, g_2, d$ & 3 \\
\midrule
\textbf{Total} & & \textbf{16} \\
\bottomrule
\end{tabular}
\end{table}

\begin{equation}
N_{\rm dof} = 13 - 16 = -3 .
\end{equation}

\section{Conclusions}
\textbf{What has been done:}
\begin{enumerate}
\item A joint quark and lepton fit with 16 parameters for 13 observables
      ($N_{\rm dof} = -3$), plus a testable axion prediction.
\item Quark sector: $\chi^2 = 6.195\times10^{-20}$.
\item Lepton sector: $\chi^2 = 1.0065\times10^{-16}$ with \(\tau_l = 0.385257 + 0.958275\,i\). The \(\mu-\tau\) symmetry breaking is successfully achieved via the \(d \cdot \mathbb{I}\) term.
\item AIC/BIC with $N_{\rm dof} < 0$ caveat.
\end{enumerate}
\textbf{What does not work:} coefficients are scale-degenerate; $\alpha_G$, $\varepsilon$, $\theta_i$ are not derived from first principles. The model interpolates observations with 16 parameters.

\appendix
\section{$A_4$ modular forms: full q-expansions}
Definition: $q = e^{2\pi i\tau}$.
\begin{align}
Y_1^{(2)} &= 1 + \sum_{n=1}^{17} C_1^{(n)} q^n, \\
Y_2^{(2)} &= -6\, q^{1/3}\left(1 + \sum_{n=1}^{11} C_2^{(n)} q^n\right), \\
Y_3^{(2)} &= -18\, q^{2/3}\left(1 + \sum_{n=1}^{10} C_3^{(n)} q^n\right),
\end{align}
$C_1 = [12, 36, 12, 84, 72, 36, 96, 180, 12, 216, 144, 36, 156, 96, 84, 240, 36]$;
$C_2 = [1, 7, 8, 18, 14, 30, 20, 42, 30, 56, 42]$;
$C_3 = [1, 2, 5, 6, 12, 8, 14, 10, 22, 16]$.

\textbf{Convolutions:}
\begin{align}
Y_1^{(4)} &= (Y_1^{(2)})^2 - Y_2^{(2)} Y_3^{(2)}, \\
Y_2^{(4)} &= (Y_3^{(2)})^2 - Y_1^{(2)} Y_2^{(2)}, \\
Y_3^{(4)} &= (Y_2^{(2)})^2 - Y_1^{(2)} Y_3^{(2)},
\end{align}
\begin{align}
Y_1^{(6)} &= \tfrac13\left[2 Y_1^{(2)} Y_1^{(4)} - Y_2^{(2)} Y_3^{(4)} - Y_3^{(2)} Y_2^{(4)}\right], \\
Y_2^{(6)} &= \tfrac13\left[2 Y_3^{(2)} Y_3^{(4)} - Y_1^{(2)} Y_2^{(4)} - Y_2^{(2)} Y_1^{(4)}\right], \\
Y_3^{(6)} &= \tfrac13\left[2 Y_2^{(2)} Y_2^{(4)} - Y_3^{(2)} Y_1^{(4)} - Y_1^{(2)} Y_3^{(4)}\right].
\end{align}

\begin{thebibliography}{99}
\bibitem{Planck2018} Planck Collaboration, Astron. Astrophys. \textbf{641}, A6 (2020).
\bibitem{PDG2024} R.~L. Workman et al. (PDG), Prog. Theor. Exp. Phys. \textbf{2022}, 083C01 (2022).
\bibitem{ADMX2025} ADMX Collaboration, Phys. Rev. Lett. \textbf{134}, 021002 (2025).
\bibitem{Feruglio2017} F.~Feruglio, Eur. Phys. J. C \textbf{77}, 713 (2017).
\bibitem{Kobayashi2019} T.~Kobayashi et al., Prog. Theor. Exp. Phys. \textbf{2020}, 053B05 (2020).
\end{thebibliography}

\end{document}